\documentclass[12pt,a4paper]{article}
\usepackage[utf8]{inputenc}
\usepackage[T1]{fontenc}
\usepackage{amsmath,amsfonts,amssymb}
\usepackage{graphicx}
\usepackage{cite}
\usepackage{url}
\usepackage{geometry}
\usepackage{fancyhdr}
\usepackage{setspace}
\usepackage{booktabs}
\usepackage{algorithm}
\usepackage{algorithmic}
\usepackage{xcolor}

% Page setup
\geometry{margin=1in}
\doublespacing

% Captions are set single-spaced and small. Under \doublespacing a 20-line caption
% doubles in height, which alone can exceed the space a float may occupy, forcing
% LaTeX to defer every figure to the end of the document.
\usepackage[font={small,stretch=1.0},labelfont=bf,labelsep=period]{caption}

% Float placement. The defaults (topfraction 0.7, textfraction 0.2) forbid a float
% taller than ~70% of the text block from sitting on a page with text, so full-width
% figures with long captions accumulate and flush at the end. These values let a
% figure appear on the page that cites it.
\renewcommand{\topfraction}{0.92}
\renewcommand{\bottomfraction}{0.90}
\renewcommand{\textfraction}{0.06}
\renewcommand{\floatpagefraction}{0.75}
\setcounter{topnumber}{2}
\setcounter{bottomnumber}{2}
\setcounter{totalnumber}{3}
\pagestyle{fancy}
\fancyhf{}
\rhead{\thepage}
\lhead{Online Social-Compliance Monitoring for AVs}

% Helper: planned content not yet supported by results (detailed protocol prose)
\newcommand{\planned}[1]{{\color{blue!60!black}\textit{[PLANNED: #1]}}}
% Claim-level pending markers (mirror claims_register.md; submission-mode CI must
% fail if any remain open). EP = internal evidence pending; XP = external pending.
\newcommand{\evidencepending}[1]{{\color{orange!75!black}\textsuperscript{[EP:#1]}}}
% \targetnum: FORMER placeholder wrapper for the C7 human-arm digits (claims
% register C7). All 14 wrapped values were replaced by the measured, independently
% verified values and unwrapped on 2026-08-20 on PI instruction (decision:
% research repo reports/knowledge/RQ022_matched_scenario_human_arm/decision.md).
% Zero \targetnum calls remain. The inert macro definition stays as a tripwire:
% any new \targetnum use marks an unapproved digit and must fail submission-mode CI.
\newcommand{\targetnum}[1]{#1}
\newcommand{\externalpending}[1]{{\color{red!65!black}\textsuperscript{[XP:#1]}}}
% Helper: figure placeholder box (keeps the file compilable without image files)
\newcommand{\figph}[2]{\fbox{\parbox[c][#1][c]{0.92\linewidth}{\centering\itshape #2}}}

% Title and author information
\title{Online monitoring of socially compliant autonomous driving}
% Author list, affiliations and corresponding-author designation supplied by the PI
% (2026-08-11); name spellings verified against jiang2025interhub. Do not alter
% without a further PI instruction.
\author{
    Xiaocong Zhao$^{1}$, Jian Sun$^{1,*}$, Xiyan Jiang$^{1}$, Yiru Liu$^{1}$,
    Gustav Markkula$^{2}$ \\[0.5em]
    \textit{$^{1}$College of Transportation Engineering, Tongji University, Shanghai, China} \\
    \textit{$^{2}$Institute for Transport Studies, University of Leeds, Leeds, United Kingdom} \\[0.4em]
    \textit{$^{*}$Corresponding author}
}
\date{\today}

\begin{document}

\maketitle

\begin{abstract}
\noindent
Autonomous vehicles are assessed for collision safety, yet crash-free interactions can depart from
human behaviour. We monitor conditional social atypicality by comparing a published trajectory-derived
interaction-preference reading with a situation-conditioned human range, withholding judgement when
the reading or human support is unavailable. Human interactions exhibit complementary roles;
conditioning narrows the reference by about one-fifth while maintaining near-nominal coverage. In
matched scenarios, automated systems are flagged about twice as often as human drivers, with the
larger excess on the accommodating side. Assertive-side flags accompany tighter routine interactions
for both vehicles without a detected increase in emergency-tail rates. A 40-participant evaluation
of selected segments finds higher counterpart-position atypicality ratings for both outside-range
categories than for within-range segments. Comfort and reported interaction cost distinguish their
directions, with larger differences on the assertive side. The framework connects departures to
subjective experience and vehicle-motion consequences while keeping social appropriateness and
realised harm as distinct questions.
\end{abstract}

% ======================================================================
\section{Introduction}
% ----------------------------------------------------------------------
Autonomous vehicles (AVs) must share the road with human drivers, cyclists and pedestrians whose
behaviour is shaped not only by physical constraints but also by implicit social norms~\cite{markkula2020defining, rasouli2020autonomous, wang2022social}. A vehicle
can be \emph{collision-safe} yet not \emph{interaction-benign}: a manoeuvre that is collision-free
in isolation may still be aggressive, hesitant or norm-violating in its interaction context,
eroding the predictability on which mixed human--machine traffic depends~\cite{trautman2010unfreezing, favaro2017examining, petrovic2020traffic}. Machine behaviour in human collectives is now an object of empirical study in its own
right~\cite{rahwan2019machine}: controlled experiments show that machine agents embedded in
human groups reshape coordination and cooperation~\cite{shirado2017locally,
crandall2018cooperating, ishowo2019transparency}, and can suppress beneficial social norms such as
reciprocity, with collective consequences~\cite{shirado2023reciprocity}. Conventional safety
verification asks ``will this trajectory collide''; a socially aware AV also faces a different
question: ``is this way of interacting still understandable and acceptable to the humans around
it?''

Two established families of methods address adjacent but distinct questions. Offline behavioural
evaluation asks whether an AV, on average, drove like a human, typically by reporting a single
aggregate sociality score over a dataset~\cite{caesar2021nuplan, montali2023waymo, sagberg2015review,
francis2025principles}, and no test-mileage budget can in any case certify rare interactive behaviour
by aggregate demonstration~\cite{kalra2016driving, koopman2016challenges}. Formal safety verification
asks whether a trajectory will collide or violate a hard constraint, and can guarantee a formally
specified notion of legal safety --- within the assumptions of its model --- as an online layer
around a planner~\cite{pek2020online, shalevshwartz2017rss, althoff2014online,
ames2019control}, an assurance now being codified in standards~\cite{iso21448sotif, ul4600}. Neither
answers the moment-to-moment question: given the current interaction state, is this behaviour within
the range that human drivers exhibit in the same situation? That question is local in time, conditional on context, and
graded rather than binary.

We make autonomous-driving social deviation \emph{visible online} by comparing current interaction
behaviour with a human reference range, while separating moments that support a judgement from those
that do not. The interpretation of a departure then requires evidence beyond interval membership.
Do people also perceive the behaviour as atypical? Do more assertive and more accommodating
departures correspond to the same comfort and interaction cost? We address these questions through
two complementary dimensions: subjective experience of selected interactions and vehicle-motion
consequences following flagged moments. Neither dimension is treated as a necessary intermediate
step to the other; direct judgements of appropriateness and evidence of realised harm remain
separate validation targets.

The monitored quantity is the interaction preference value (IPV), a trajectory-derived measure of
how an agent balances its own progress against the group's. The trajectory-based estimator was
introduced and validated in controlled VGIM interactions in our previous work~\cite{zhao2024measuring};
here we use its frozen online implementation as the behavioural reading supplied to the
human-reference monitor. A related inverse-game formulation provides complementary evidence for
quantifying social patterns from trajectories~\cite{liu2024sociality}.

Three distinctions bound the claims. First, \emph{socially atypical is not socially inappropriate}:
lying outside the human range is a statement about typicality. Subjective experience and
vehicle-motion outcomes help interpret that departure without converting it into a normative
verdict. Second, \emph{collision-safe is not interaction-benign}. Third, \emph{not readable is not neutral}: before
an interaction is behaviourally informative or after it has resolved, the IPV may be weakly
identifiable or not applicable, and this must not be read as compliant. A responsible social
monitor must know when \emph{not} to judge; abstention is an explicit boundary of responsible
judgement, not a failure mode, in the tradition of prediction with a reject
option~\cite{chow1970optimum, cortes2016learning, geifman2017selective}. The human range is, moreover,
a \emph{reference range} in the sense familiar from clinical medicine and normative
modelling~\cite{ozarda2016reference, marquand2016understanding, bethlehem2022brain}, not a demand to
imitate one human trajectory.

This framing differs from three existing paradigms. Unlike offline AV
evaluation~\cite{caesar2021nuplan, montali2023waymo}, it is online, per-timestep and
context-conditioned. Unlike runtime verification of formal
specifications~\cite{leucker2009brief, maler2004monitoring, sha2001using, lindemann2023conformal}, the
monitored property is an empirical reference distribution estimated from human data. Unlike socially
compliant planning~\cite{schwarting2019social, sun2018courteous, sadigh2016planning, kruse2013human,
mavrogiannis2023core}, it does not encode sociality into a planner's objective but provides a
pluggable, planner-agnostic monitor that any planner can be checked against. Closest to our aim,
online monitors exist for codified legal rules~\cite{yu2024online}; unusual driving can be flagged as
unexplained deviation from learned traffic behaviour~\cite{hu2023socially} or as violated
expectation~\cite{dinparastdjadid2023measuring, engstrom2024modeling}; and AV conduct is benchmarked
offline against reference models of a competent human driver~\cite{scanlon2026collision,
schumann2026active}. None of these reads a social preference, refers it to a calibrated
situation-conditioned human range, or separates the moments at which a social judgement is meaningful
from those at which it is not.

This work reframes social compliance as the online monitoring of conditional social atypicality.
Human interaction is relational and context-dependent: partners take complementary roles, yet the
reference needed online can be determined from the current observable situation without reading a
hidden intention. Built on that situation-conditioned range, the monitor is calibrated and auditable
at runtime, and it abstains whenever the interaction is unreadable or the situation lacks comparable
human support. We then apply the unchanged instrument to a matched-scenario real-vehicle benchmark,
where 19 automated systems and 20 licensed drivers face the same 15 interactions. Automated systems
are flagged about twice as often as the matched human drivers, with the larger excess on the
accommodating side. At assertive-side flagged moments, ordinary interaction quality tightens for both
vehicles without a corresponding increase in emergency-tail events. The human arm checks alarm
inflation after transfer to a new country and apparatus and examines the same consequence signature
in human-driven interactions. A separate subjective evaluation with 40 participants asks how
90 selected interaction segments are experienced from the ego and counterpart positions. It links
both outside-range categories to perceived atypicality from the counterpart position and identifies
larger comfort and reported-cost differences on the assertive side. Together, the studies separate
how frequently a departure occurs, its direction, and its experiential and motion-related meaning.

% ======================================================================
\section{Results}
% ----------------------------------------------------------------------

\subsection{A human-reference monitor makes conditional social atypicality visible online}

The trajectory-based IPV estimator was previously validated in controlled VGIM
interactions~\cite{zhao2024measuring}; here it is held fixed and used as the frame-level behavioural
reading supplied to the human-reference monitor. Figure~\ref{fig:measured} summarises the monitoring
pipeline, including the candidate-based reading, the context-conditioned reference, explicit
abstention and a complete interaction example.

\begin{figure}[tbp]\centering
\includegraphics[width=\linewidth]{figures/fig1_three_layer.pdf}
\caption{\textbf{The human-reference monitoring framework, from a trajectory-derived IPV reading to a context-conditioned verdict.}
\textbf{a}~The reading is an angle: a candidate preference $\theta$ weights the agent's own-progress
cost by $\cos\theta$ and the interaction cost by $\sin\theta$, so $\theta$ is the direction in which
the two are traded. Seven candidates span $[-3\pi/8,\,3\pi/8]$ at $\pi/8$ spacing; below the
horizontal axis the counterpart's cost carries negative weight, the assertive side.
\textbf{b}~Two frames of one interaction, five apart. Bars are the normalised candidate likelihoods
of Eq.~\eqref{eq:estimator}: concentrated above, where the reading is their weighted mean; at the
uniform value $1/7$ below, where the frame does not say which candidate the agent is following and
the monitor abstains.
\textbf{c}~Human reference readings on the held-out fold ($486{,}660$ moments) and the range applied
to them. Bars give the median $90\%$ range in three real situation classes, narrowest to widest:
crossing with priority, merging with priority, merging with equal priority; shading gives the verdict
the middle class returns. The deployed range also conditions on the continuous situation, so it
varies within a class ($1.05$--$2.37$\,rad wide across the benchmark; Methods).
\textbf{d}~The same population through both gates, from $663{,}282$ candidate moments: $486{,}660$
readable, $461{,}937$ of those also human-supported, splitting $417{,}036$ inside and $44{,}901$
outside. Failing either gate returns no verdict, which is not a pass.
\textbf{e}~Local trajectories; time markers link to \textbf{f} and the red circle is the conflict
zone. \textbf{f}~The same interaction frame by frame, drawn only where both gates hold, so grey
intervals denote no reading rather than a zero. The three ranges move with the situation; the marked
readings lie above the $90\%$ range, on the accommodating side. Panels \textbf{e} and \textbf{f} are
descriptive. Source data are provided as a Source Data file.}
\label{fig:measured}
\end{figure}

\subsection{Human interaction runs on a two-sided script, and the reference must condition on the situation}

We next ask what a reading should be compared against, beginning with the structure of human
interaction in the space that the reading describes. We analysed $38{,}228$ interaction cases from
four naturalistic sources within the InterHub corpus~\cite{jiang2025interhub} (Waymo $23{,}218$;
nuPlan $7{,}499$; Lyft $5{,}105$; Argoverse-2 $2{,}406$; $3{,}695{,}981$ trajectory frames),
summarising each agent by the mean IPV over its valid interaction segment. Seen whole, an interaction
between two humans follows a two-sided script rather than two independent scores. Plotting each valid
case as the pair of its partners' summaries concentrates the mass where one partner sits below zero
while the other sits above it. Across the $34{,}757$ cases in strata that
support matched re-pairing, the within-interaction correlation is $-0.338$ (95\% CI
$[-0.351,-0.325]$), whereas re-pairing partners across cases within matched scenario strata gives
$+0.041$ ($[+0.030,+0.052]$); an independently built implementation returns $-0.339$, and the
correlation is negative in every source.

The division of roles is legible early. Readings from the first $5\%$ of an interaction identify
which partner will ultimately give way at AUC $0.75$, rising to $0.81$ by half-way, while early
kinematics over the same windows remains near $0.56$--$0.57$. The script is
also keyed to the situation before the conflict is joined: turning vehicles read $+0.079$\,rad
further toward giving way than their straight-driving partners ($[+0.063,+0.095]$,
$n=3{,}565$ pairs), in the same direction in all four sources, with one source's interval admitting
zero. Social behaviour at this level is relational---organised by partner,
role, geometry and risk---rather than a fixed disposition of an individual.

No single number attached to a role reproduces that script. What holding right-of-way contributes is
not a fixed offset: across three bands of realised post-encroachment margin the priority minus
non-priority difference is $+0.058$ ($[+0.050,+0.067]$, $n=9{,}799$), $+0.001$
($n=22{,}833$) and $-0.034$ ($[-0.045,-0.023]$, $n=4{,}863$; Fig.~\ref{fig:context}a).
These offline bands are used only to describe the human data and are never read by the online monitor.
A coarse geometric prior points in the same direction across all four sources
(Fig.~\ref{fig:context}b), and the dependence is preserved
under dropping the largest source and alternative geometry and risk binnings.

The cost of a single global reference is visible in the distribution itself. Across the $486{,}660$
readable human moments of the held-out fold, the range selected by the situation is narrow in some
real situations and wide in others: among three illustrative classes, crossing with priority is the
narrowest and merging with equal priority the widest, while the deployed range spans
$1.05$--$2.37$\,rad across the benchmark (Fig.~\ref{fig:measured}c). The classes make the variation
visible; the deployed range conditions on the continuous situation and varies within each class. We
report this as a regularity of the observed corpora rather than a transferable law: fitting the full
state description on every source but one explains at most $+0.03$ of the variance in the held-out
source and falls below zero in two (Fig.~\ref{fig:context}c; Methods 4.4).

% Interim caption matches the existing three-panel asset inspected in the base PDF.
% The local figure pass must restore the six-panel semantic evidence display and
% update its caption and panel callouts together; see the figure handoff plan.
\begin{figure}[tbp]\centering
\includegraphics[width=\linewidth]{figures/fig2_context.pdf}
\caption{\textbf{Human interaction depends on role and geometry, with a boundary on cross-source prediction.}
\textbf{a} Priority-minus-non-priority mean IPV in three realised post-encroachment-time bands.
The difference changes from $+0.058$ to $+0.001$ to $-0.034$\,rad across the displayed bands;
these are descriptive offline strata, not inputs to the online monitor.
\textbf{b} Coarse geometry effects on the reading in each of four sources: merge-or-pass and
same-direction pairings show directionally consistent shifts with source-dependent magnitudes.
\textbf{c} Variance explained when the state-based predictor is fitted without the evaluated
source. Values near or below zero delimit transfer of the state-to-reading relationship.
Error bars show $95\%$ confidence intervals where available. The complementary-role, early-legibility
and turning-partner results are reported in the text and Methods; this existing asset displays the
role, geometry and source-transfer analyses. Source data are provided as a Source Data file.}
\label{fig:context}
\end{figure}

\subsection{A situation-conditioned human range is a calibrated monitor}

Establishing that human behaviour is context-dependent does not by itself say what an online monitor
must read. We define online monitoring as a membership test: the ego's current IPV should lie within
a human reference range determined by the current observable situation,
\begin{equation}
\begin{gathered}
  \hat\theta_i(t)\ \stackrel{?}{\in}\ I_{\mathrm{human}}(z_t),\\
  z_t=\big(\text{geometry, role, progress, online-computable risk proxies,}\\
  \text{relative kinematics, support}\big),
\end{gathered}
  \label{eq:membership}
\end{equation}
evaluated only where the interaction is active and the ego is readable; the risk proxies are those
computable from information available at time $t$ (no observed post-encroachment measure enters
online). The range is the conditional
distribution of the human current IPV given $z_t$, calibrated by split-conformal
prediction~\cite{vovk2005algorithmic,romano2019cqr} and
abstaining when the state is unsupported. Conditioned on the current situation, the human range is
sharper than a single global range---about a fifth narrower at the 80\% and 90\% levels (21\% and
20\%), at coverage within 0.6 percentage points of nominal---and it stays informative at the
strictest level, where the global range spans essentially the whole admissible scale and retains
almost no power to discriminate. In absolute terms the conditioned $90\%$ range still averages $1.87$\,rad
(5th--95th percentiles $1.35$--$2.28$\,rad) of the $2.36$\,rad candidate span: the reference is wide,
and a flag therefore marks a clear departure from the human range rather than a fine discrimination.
The coverage statement is empirical and marginal over accepted moments; it is not a
per-situation or per-interaction guarantee (Methods 4.4). Where a reading is available, the monitor abstains for lack of human support on only a small
fraction of readable moments ($5.08\%$ on the held-out natural-driving test fold; abstention on the real-vehicle
benchmark is larger and is reported with its reasons in Section 2.4 and Methods), giving an auditable runtime monitor
of conditional social atypicality (Fig.~\ref{fig:monitor}). The range is not a fixed corridor:
it rises, falls and changes width as the situation changes, so that the same reading can sit inside it
in one situation and outside it in another (Fig.~\ref{fig:edband}).

\begin{figure}[tbp]\centering
\includegraphics[width=\linewidth]{figures/fig3_monitor.pdf}
\caption{\textbf{The situation constrains social preference without determining it, which is why the
reference range must be conditioned and the preference itself must be read online.} Grey marks a single
global human range; teal marks the range conditioned on the current observable situation
(Eq.~\ref{eq:membership}).
\textbf{a} Width of the reference range at three nominal coverage levels, as a percentage of the
admissible range of the quantity. The admissible range is the span of the candidate grid,
$3\pi/4\approx2.36$\,rad; the parameterisation domain of Eq.~\eqref{eq:estimator} is $[-\pi/2,\pi/2]$ and the seven
candidates span $[-3\pi/8,3\pi/8]$. A global range is not a usable reference: at the $95\%$ level its
width equals $100.00\%$ of the admissible span, so it is as wide as the quantity itself and retains
almost no power to discriminate --- it still flags, but only where the reading sits at the very edge
of what the estimator can express, and it pays for that width by over-covering (panel~\textbf{b}).
Conditioning narrows it by
$21\%$, $20\%$ and $8\%$ at the three levels (reductions computed from unrounded widths --- $20.7\%$,
$20.5\%$, $8.4\%$; the narrowing labels drawn on the panel are rounded to whole percent, so
recomputing from them can differ by a point) (at the $90\%$ level the conditioned width has mean
$1.87$ rad, 5th--95th percentiles $1.35$--$2.28$ rad).
\textbf{b} Achieved minus nominal coverage. The conditioned range sits within $0.6$ percentage points
of nominal at all three levels ($+0.03$, $+0.28$, $+0.57$); the global range over-covers by $2.8$
points at $95\%$.
\textbf{c} Variance of the online reading that the situation accounts for, out of fold. The situation
explains $20.9\%$ (defined as $1-\mathrm{SSE}/\mathrm{SST}$ of the conditional median against the
test-fold mean, over all $486{,}660$ test moments; $21.1\%$ over the $461{,}937$ accepted moments); the remaining $79.1\%$ is not recovered
by the situation description---which is what makes an online reading necessary rather than
redundant.
Panels \textbf{a},\textbf{b} use $n=461{,}937$ evaluated moments (the accepted moments; identical at
all three levels); \textbf{c} uses $n=486{,}660$. The
monitor abstains on $5.08\%$ of moments, where the situation falls outside the support of the human
record. It is empirical: it takes no observed risk or post-encroachment measure as an online input, and
it returns an explicit abstention rather than a neutral verdict where the situation is unsupported.
Verdicts elsewhere in this paper are issued at the $90\%$ level.
Transfer to a data source not seen during fitting is not established and is reported as a boundary of
the present monitor. Source data are provided as a Source Data file.}
\label{fig:monitor}
\end{figure}

A striking feature organises this result. Although the interacting partners' IPVs are correlated at
the event level (above), \emph{adding the counterpart's own inferred preference to the
situation does not measurably sharpen the reference range}; conditioning on the ego agent's own
behavioural history does not either. Both are retained only as ablations. Thus the event-level
complementarity does not translate into measurable additional sharpening from the counterpart-IPV
channel in this implementation. The current observable situation supplies the deployed reference
without requiring that additional inferred preference. The monitor returns an explicit abstention
when the situation lacks human support or the reading is unavailable. Its supported interpretation
remains conditional social \emph{atypicality}.

\subsection{From atypicality to interaction consequences}
We ask what an atypicality verdict corresponds to in the interaction that follows, when
the same situations are driven by different automated systems. We pursue this on a matched-scenario
real-vehicle benchmark~\cite{onsite}, in which identical scenarios are replayed against many
independent driving systems.

The monitor is, first, able to speak. The benchmark replays 15 scenario templates against 19
independent driving systems (267 scenario runs). Across its 67{,}861 candidate moments the monitor
returns a verdict on 14{,}099 ($20.8\%$), with no engineering failures in this replay --- no solver or pipeline failure on any
candidate moment --- and in 231 of the 267 scenario runs ($86.5\%$) it speaks at least once. At the
moment level, $55.3\%$ of candidate moments carry discriminative information between the candidates,
and $20.8\%$ additionally lie within human support; at the run level, none of the 36 runs in which the
monitor never speaks is silenced by readability alone --- in every such run there are readable
moments, and what is missing there is comparable human behaviour for the situation. Abstention here
is a statement about the reference rather than about the vehicle, and it is auditable as such.

If human interaction follows a two-sided script (Section 2.2), a departure by one partner should be
visible both in that vehicle's own margin and in what the counterpart absorbs. Where the monitor does
speak, atypicality on the assertive side---a current preference that weights the counterpart's cost less
than humans do in the same situation---is followed by a measurably tighter interaction for
\emph{both vehicles}. The ego's own margin to its counterpart contracts at the
comfortable end of its distribution: the upper quartile falls by two-fifths, while the lower
quartile---the end where a margin would actually be pressing---is not compressed at all. The
counterpart, read independently from the other vehicle's own control record, carries nearly
twice the routine speed variation over the same window, and roughly twice the routine speed reduction
($+2.6$ and $+1.4$\,km\,h$^{-1}$ as differences, which do not depend on how small the routine value
is). The two
sides are measured from different sources and different quantities, so their agreement is a
convergence rather than a restatement (Fig.~\ref{fig:consequence}); the association survives a
placebo test that reassigns whole flag sequences across scenario runs (Methods 4.5); a single run containing flagged
assertive moments is shown for illustration in Fig.~\ref{fig:edcase}.

The same outcome measures, applied unchanged to the accommodating side, do not reproduce that
signature: neither the counterpart's speed variation nor its speed reduction separates detectably
from within-range moments, and neither the median nor the upper quartile of the ego's own margin
does either. The two sides are also compared
with each other directly, rather than each against the within-range group in turn, and on the
quartile where the assertive side compresses they differ (Methods 4.5).
The result is direction-specific: assertive-side departures accompany compression of the measured
routine vehicle-motion outcomes, whereas accommodating-side departures show no corresponding
compression in this battery. The latter finding concerns these objective endpoints; subjective
experience is examined separately in Section~\ref{sec:subjective_results}.

% Existing asset retained in the text-only pass. Its counterpart-response forest
% and threshold-difference intervals show the assertive comparison; the local
% figure pass must add the published accommodating comparisons without new tests.
\begin{figure}[tbp]\centering
\captionsetup{singlelinecheck=false}
\includegraphics[width=\linewidth]{figures/fig5_consequence.pdf}
\caption{\textbf{Assertive-side flags accompany tighter routine interactions without a detected increase in emergency-tail events.}
Rows separate the ego vehicle from the counterpart; columns separate routine outcomes from emergency
thresholds. Teal denotes within-range moments, dark red the assertive side and light red the
accommodating side. The three categories appear in the margin distribution and threshold-share
summaries; the interval summaries in this existing display show assertive-versus-within contrasts.
The accommodating-side intervals are reported in Methods 4.5.
\textbf{a} Cumulative distribution of the minimum ego-to-counterpart time-to-collision over the
three-second post-verdict window, omitting non-closing frames. The lower quartile is not compressed
($3.46$ versus $2.99$\,s); the upper quartile falls $39.9\%$ $[16.9,56.4]$, while the median contrast
is unresolved ($9.9\%$ $[-10.7,23.7]$). Direct assertive-minus-accommodating comparisons distinguish
the upper-quartile ratios ($-0.43$ $[-0.71,-0.06]$) but not the median ratios
($-0.09$ $[-0.29,+0.12]$).
\textbf{b} Shares below all four prespecified margin thresholds, with assertive-minus-within
share differences and intervals. Supported differences point toward fewer short margins.
\textbf{c} Counterpart-response median ratios relative to within-range moments. Assertive-side
speed reduction is $2.74$ versus $1.33$\,km\,h$^{-1}$ and speed range is $5.53$ versus
$2.93$\,km\,h$^{-1}$. Corresponding median differences are $+1.41$ $[+0.08,+3.37]$ and
$+2.60$ $[+1.01,+3.50]$\,km\,h$^{-1}$. Grey steering entries have intervals admitting parity.
\textbf{d} Counterpart-record shares below each braking threshold, with assertive-minus-within
differences and intervals; this is the only record-weighted panel.
Panels \textbf{a},\textbf{b} contain $499$ assertive, $817$ accommodating and $12{,}344$ within-range
moments; \textbf{c},\textbf{d} use $469$, $719$ and $10{,}483$ anchors with counterpart logs.
Panel \textbf{d} pools $13{,}800$, $21{,}025$ and $310{,}246$ acceleration records, respectively.
Intervals are pointwise $95\%$ bootstrap intervals over whole scenario runs: $227$ ego-window runs
and $175$ counterpart-log runs. Quantile ratios and speed contrasts use $2{,}000$ draws; threshold
contrasts use $1{,}000$. Thresholds are nested. The unchanged outcome battery and its full
side-specific results are described in Methods 4.5. Associations are descriptive. Source data are
provided as a Source Data file.}
\label{fig:consequence}
\end{figure}

Counting both sides, $9.8\%$ of judgeable moments (candidate moments that pass both gates; Methods 4.4)
are flagged ($1{,}388$ of $14{,}099$; $519$ on the assertive side and $869$ on the accommodating side).
These fractions describe the frequency and direction of departures. The outcome analyses describe
what follows flagged moments. The next section places the flag rate alongside a matched human control.

This is a compression of ordinary interaction quality, not a shift of mass into the tail. At the
emergency thresholds---a closing margin short enough to demand immediate action, or braking hard
enough to count as evasive---assertive-side moments are roughly half as frequent as typical ones
($0.46$ to $0.61$ times) at every threshold whose interval excludes a zero difference (margins under
$1$, $1.5$ and $2$\,s; braking beyond $-2$, $-3$ and $-4$\,m\,s$^{-2}$); at the widest margin tested
(under $3$\,s) the interval admits no difference. On the accommodating side the only threshold whose
interval excludes zero is the tightest margin, and it points the same way, with fewer short margins
rather than more. The lower quartile of the ego's margin is not compressed on either side, sitting
marginally wider for flagged moments than for typical ones. What
the monitor registers is that ordinarily comfortable margins tighten, not that dangerous events become more common.
The accommodating side shows no corresponding compression of these routine vehicle-motion outcomes.

Compression of interaction quality is not yet a harm claim, and we do not make one. Under the
present definitions of a realised interaction failure, atypicality does not robustly or specifically
predict harm outcomes; the binding limitation is upstream, in that identifying and segmenting
interaction-failure events is itself not yet reliable. The present evidence therefore bounds the
harm relationship rather than confirming it.

What, then, does a social flag register that a conventional safety check does not? The routine
interaction differences above occur without an observed rise in the tested emergency-tail rates.
Thus the social reading and the threshold-based safety endpoints describe different aspects of
these interactions; the routine differences are not evidence of more emergency events.

\subsection{The reference is audited by the population that defined it}
% C7: the human-arm numeric literals in this section are the measured,
% independently verified values (swapped and unwrapped 2026-08-20; decision:
% research repo reports/knowledge/RQ022_matched_scenario_human_arm/decision.md).
The reference range was learned from natural human driving recorded in other countries on open
roads, and every verdict above rests on the premise that this reference remains readable in a new
country, on a closed course, in deliberately staged conflicts. We test that premise directly:
licensed human drivers drove the same 15 scenarios in a real vehicle, against the same
counterparts under the same control, and the unchanged monitor---same estimator, same gates, same frozen
reference---audits them exactly as it audits the automated systems.

Three observations follow (Fig.~\ref{fig:humanarm}). First, the instrument survives the move. The
population that defined the reference lies outside it on $9.72\%$ of held-out judgeable moments at
the $90\%$ level---the instrument's native alarm rate. Human drivers on the course are flagged on
5.0\% of judgeable moments (786/15{,}598), no more often than
that native rate: carrying the reference across country, apparatus and task does not inflate its
alarms. The course human rate therefore serves as the matched calibration control; the automated systems are
evaluated against the same frozen reference in the same scenarios, and their corresponding rate is
reported below. The two settings are not, however, interchangeable. Drivers on the course were asked to drive
as they normally would, but the scenario set is not ordinary driving: these are staged conflicts
selected for a benchmark, and they are harder than what a naturalistic corpus mostly contains. The
cross-corpus reading is therefore an alarm-inflation check---does the instrument fire more often
once it is carried somewhere new---and not a comparison of levels. The comparison of levels is the
one made inside the course, where both kinds of driver meet the same fifteen scenarios and the same
counterparts under the same control. How often the monitor fires at all is a property of the
scenarios it is pointed at: across these fifteen the human flag rate runs from $1.7\%$ to $12.4\%$
and the automated rate from $2.7\%$ to $17.4\%$, so an absolute alarm rate is not a portable
quantity. That is why the audit carries its own human arm rather than treating the corpus rate as a
yardstick. The yardstick itself never changes: both arms are judged against the same frozen
reference, built from human--human interaction alone, and the question asked of each is whether the
behaviour falls inside the human range for this situation---never whether it resembles the other
arm. Second, the flag means the same thing for either kind of
driver. At the moments where a human is flagged on the assertive side, the same consequence signature appears that
Fig.~\ref{fig:consequence} documents for automated vehicles: the ego margin's upper quartile
contracts by about three-tenths, the counterpart absorbs $2.5$ times the routine speed
reduction, and every emergency tail is rarer, not more common. The verdict marks a property of the
interaction, not a property of machines. Third, under this common instrument the automated systems
are flagged $2.0$ times as often as the matched human drivers ($1.95$, $95\%$ CI $[1.62, 2.39]$
resampling scenarios; $9.8\%$ vs $5.0\%$ of judgeable moments), and the per-scenario flag rate is
higher for the automated systems in 15 of 15 scenarios---in every resampled scenario set the ratio
stays above parity. The population-level rates answer
a calibration question; the moment-level signature answers what a flag means; neither is a ranking
of vehicles or drivers.

The side-specific split sharpens the comparison. Automated systems exceed matched humans primarily
on the accommodating side ($6.16\%$ versus $2.25\%$), whereas the assertive-side difference is
smaller ($3.68\%$ versus $2.79\%$). The larger machine--human excess therefore lies on the side for
which the frozen vehicle-motion battery detects no corresponding compression signature. Whether
that side also differs in human experience is a separate question, addressed next.

% Interim caption matches the existing a-c asset; the missing side-split panel
% is specified in the figure plan and its measured values remain in the text.
\begin{figure}[tbp]\centering
\includegraphics[width=\linewidth]{figures/figH_human_arm.pdf}
\caption{\textbf{The monitor is audited by the population that defined its reference.}
\textbf{Top}~Moments entering the comparison in each arm: candidate moments, those that are
readable, and those additionally judgeable against the reference; for the automated systems the
middle stage is given as the fraction passing. The human protocol is a fixed scenario sequence,
identical across drivers, with no missing trial.
\textbf{a}~Fraction of judgeable moments flagged outside the reference at each nominal level, for
the held-out natural-driving reference population, matched human drivers and automated systems;
dashes mark nominal outside rates. At $90\%$, the human flagged moments comprise $435$ assertive
and $351$ accommodating moments, and the automated flagged moments comprise $519$ and $869$.
The human same-scenario $90\%$ outside rate has a driver-by-scenario bootstrap $95\%$ interval
of $3.4$--$5.5\%$; displayed rates are point estimates.
\textbf{b}~Consequence signature at flagged-assertive moments, expressed as flagged-to-within
ratios; the braking row counts counterpart records below $-3$\,m\,s$^{-2}$. Filled markers denote
automated vehicles and open markers human drivers. Asterisks denote case-bootstrap intervals
excluding parity, or excluding zero share difference for the emergency endpoints. Whiskers appear
where intervals were computed on the ratio. The two human ego-margin entries have no interval in
the frozen battery; absence of an interval is not a null result.
\textbf{c}~Scenario-specific human and automated flag rates at $90\%$; the diagonal marks parity
and the 15 scenario points are shown without intervals. Both arms use the same frozen human--human
reference; no panel ranks individual systems or drivers. Source data are provided as a Source Data file.}
\label{fig:humanarm}
\end{figure}

\subsection{Subjective experience distinguishes the two directions of atypicality}
\label{sec:subjective_results}
% Adopted for exploratory manuscript reporting on PI instruction, 2026-09-25.
% Source: ipv_estimation@0e05a65, final RQ028 extraction package.
We examined whether the monitoring categories were also reflected in human experience of selected
interaction segments. This separate study comprised 40 participants in 20 pairs and 90 segments
from 15 scenarios; 596 completed trials yielded 1{,}192 position-specific rating records. From the
counterpart position, both assertive-side and accommodating-side segments were rated as more
atypical than within-range comparison segments: mean within-participant differences were
$+0.639$ $[+0.318,+0.992]$ and $+0.682$ $[+0.548,+0.817]$ points, respectively, on a 1--7 scale.
These contrasts survived Holm correction across all 30 exploratory comparisons. Thus the supplied
monitoring categories have a perceptual counterpart in this selected stimulus set
(Fig.~\ref{fig:subjective}; Supplementary Table~\ref{tab:subjective_core}).

Comfort and reported interaction cost distinguished the two directions more clearly. Relative to
within-range segments, assertive-side segments received counterpart-position comfort ratings
$1.091$ points lower and interaction-cost ratings $2.070$ points higher. Direct comparison with
accommodating-side segments gave a comfort difference of $-0.870$ $[-0.974,-0.766]$ points and a
reported-cost difference of $+1.163$ $[+1.019,+1.300]$ points. Ego-position comfort followed the same
ordering, with the assertive category below both other categories. The result therefore concerns
which side of the reference is crossed as well as whether a segment is outside it.

Accommodating-side segments also differed from within-range segments: counterpart-position comfort
was $0.221$ points lower ($[-0.351,-0.100]$) and reported interaction cost was $0.908$ points higher
($[+0.771,+1.051]$). These smaller subjective differences complement the absence of a corresponding
compression signature in the separate vehicle-motion analysis. They describe another dimension of
interaction quality; the two studies do not establish a segment-matched dissociation or a causal
effect of crossing a reference boundary.

The main comfort and cost directions were retained when categories were compared within shared
scenarios. For assertive-versus-within comparisons, counterpart comfort was lower and reported cost
higher in all eight shared scenarios; the corresponding assertive-versus-accommodating directions
held in all 12 shared scenarios. These are scenario-composition sensitivity checks, not independent
replications. Estimates and intervals preserve participant pairing and condition on the selected
stimuli (Methods~\ref{sec:subjective_methods}).

% Figure rendering is intentionally deferred to the local agent. The placeholder
% marks missing artwork, not hypothetical data. Do not remove it until the asset
% and its panel-level source manifest have been checked.
\begin{figure}[tbp]\centering
\IfFileExists{figures/fig6_subjective.pdf}{%
\includegraphics[width=\linewidth]{figures/fig6_subjective.pdf}%
}{\figph{1.2in}{Figure 6 artwork pending local rendering.\\
Measured endpoints: perceived atypicality, comfort and reported counterpart interaction cost.}}
\caption{\textbf{Subjective ratings distinguish perceived atypicality from interaction experience.}
The planned display uses the measured ratings of assertive-side (A), accommodating-side (C) and
within-range comparison (W) segments. Perceived atypicality is the reversed typicality item
($8-\mathrm{Q1}$); higher Q4 ratings indicate greater comfort; Q5 is the reported counterpart
interaction-cost item. Each position contributes 158 A, 239 C and 199 W completed ratings from
40 participants in 20 pairs. Class means first average repeated ratings within participant and
then give participants equal weight. Intervals are pointwise $95\%$ intervals from 10{,}000
participant-pair bootstrap draws, conditional on the selected stimuli. Q5 is available only at the
counterpart position. Main contrasts are tabulated in Supplementary Table~\ref{tab:subjective_core};
the full 30-comparison family uses Holm correction. Final panel labels will be set with the artwork.}
\label{fig:subjective}
\end{figure}

% ======================================================================
\section{Discussion}
% ----------------------------------------------------------------------
This work makes social monitoring an auditable question about behaviour in a situation: whether the
current interaction lies within a human reference range, when that comparison is supportable, and
what an observed departure means. Human behaviour is relational and context-dependent, with
complementary roles that are legible early. The situation-conditioned reference is sharper than a
global range at near-nominal coverage and withholds judgement when the reading or human support is
missing. In matched real-vehicle scenarios, the reference does not inflate alarms for human drivers,
and assertive-side flagged moments show a shared vehicle-motion signature in human and automated
arms. The separate subjective study adds evidence that the selected outside-range segments are
perceptibly atypical and that their direction matters for comfort and reported interaction cost.

Two distinctions organise these findings. First, interaction-level complementarity and the inputs
needed by an online reference are different questions. Although partners' readings covary across
interactions, adding counterpart IPV or ego self-history does not measurably sharpen the reference
beyond the current situation in the evaluated implementation. The deployed monitor can therefore
use the situation-conditioned range without those additional channels. This ablation result does not
identify a psychological mechanism behind the partner relationship.

Second, the frequency of departures and their experiential meaning need to be reported separately.
The larger automated--human excess occurs on the accommodating side, while the larger comfort and
reported-cost differences in the subjective study occur on the assertive side. Accommodating-side
segments nevertheless carry smaller subjective differences relative to within-range segments.
Together, these observations motivate retaining direction when interpreting a flag rather than
using its frequency as a proxy for interaction cost. The rates and ratings come from different
analyses and are not combined into a population-level social-loss score.

Subjective experience and vehicle-motion consequences also provide complementary endpoints. The
real-vehicle result is compression of routine interaction quality for the ego and counterpart at
assertive-side flagged moments, without a detected increase in the tested emergency tails. The
absence of the same compression on the accommodating side applies to that fixed motion-based
battery. Its lower comfort and higher reported-cost ratings in the subjective study show why a
broader evaluation should retain an experiential dimension. The current studies do not resolve
which individual motion features account for those ratings.

The empirical reference, subjective experience, direct normative judgement and realised harm remain
distinct. The questionnaires support statements about perceived atypicality, comfort and the
reported interaction-cost item; these endpoints do not establish psychological IPV recovery or a
universal judgement of social appropriateness. The subjective analysis is exploratory and its
participant-pair intervals condition on 90 selected segments. Within-scenario comparisons check
composition sensitivity without establishing generalisation to all new stimuli. A reliable link to
realised interaction failure remains unconfirmed under the present outcome definitions.

The human-driving arm carries a separate scope: participants drove an instrumented vehicle on a
closed course against controlled counterparts, supervised and instructed to drive naturally, with a
fixed scenario sequence. Its protocol and consent cannot be substituted for those of the subjective
study. The evidence covers pairwise vehicle--vehicle interactions at mapped conflict points;
cyclists, pedestrians and multi-agent scenes remain outside the present reference. Beyond driving,
the framework offers a template for comparing autonomous-agent behaviour with human reference
ranges while separately examining experience and interaction outcomes and explicitly withholding
unsupported judgements~\cite{rahwan2019machine}.

% ======================================================================
\section{Methods}
% ----------------------------------------------------------------------
\noindent\emph{Note on scope.} The main text states the social-monitoring argument; this section
records the estimator, the reference-range model and its ablations, the calibration and abstention
machinery, the real-vehicle consequence protocol, the matched human-driving arm and the separate
subjective-evaluation analysis. Throughout, \emph{online} means that every input to a
verdict is available at decision time (non-anticipation verified); end-to-end real-time operation on
vehicle hardware is not evaluated in this paper.

\subsection{Online likelihood-based inverse-planning IPV estimation}
The trajectory-based IPV estimator was introduced and validated in controlled VGIM interactions,
where imposed interaction-preference patterns were identified from trajectories and improved
interactive-motion prediction~\cite{zhao2024measuring}. We adopt the published estimator as a fixed
component of the monitor; Supplementary Note 1 summarises the earlier validation and the present
online implementation.

For a grid of candidate preferences $\{\theta_k\}$, candidate trajectories
$\widetilde X_a(W_t;\theta_k)$ are generated and scored against the observed trajectory, the
counterpart trajectory and the static map-lane-centreline reference. The frame-level reading is the
likelihood-weighted mean,
\begin{align}
  \ell_{a,k}(t) &= p\!\left(X_a(W_t)\,\middle|\,\widetilde X_a(W_t;\theta_k),\,X_j(W_t),\,\mathcal{M}\right),
  & \pi_{a,k}(t) &= \frac{\ell_{a,k}(t)}{\sum_m \ell_{a,m}(t)},
  & \hat\theta_a(t) &= \sum_k \pi_{a,k}(t)\,\theta_k,
  \label{eq:estimator}
\end{align}
where $W_t$ is the interaction window up to time $t$, $X_j$ is the counterpart trajectory and
$\mathcal{M}$ is the map and reference geometry. The IPV is parameterised as
$\theta\in[-\pi/2,\pi/2]$, with positive values placing more weight on the interaction term.
Each candidate trajectory for preference $\theta_k$ optimises the frozen implementation's
utility $\cos\theta_k\cdot(\text{own-progress cost}) + \sin\theta_k\cdot(\text{interaction cost})$, so
a more negative $\theta$ weights the joint term negatively --- the assertive direction; the full cost
terms and solver configuration are released with the code. This is inverse planning, not a single
joint maximum-likelihood optimum; $u_a(t)$ is a dispersion
statistic of the normalised candidate weights, not a standard deviation, and the operational readability gate
$u_a(t)<u_0$ of Eq.~\eqref{eq:readability} is operationalised by the following frozen rule: the
reading is discarded as near-uniform when the largest normalised weight falls below $0.20$ (against
$1/7\approx0.14$ for exactly uniform weights) or the candidate likelihoods are exactly tied. The
reference frame is what makes the estimate usable online: against an observed-trajectory-prefix
reference the causal current estimate agrees only weakly with the offline estimate (correlation
$\approx 0.281$), whereas against a static map-lane-centreline reference, route-conditioned and
available before the outcome, it agrees closely (correlation $\approx 0.993$, MAE $\approx 0.027$),
bounded to the lane-referenced slice (about three quarters of cases) and stated as a non-anticipation
property, not a causal-effect claim. The gate determines whether the frozen candidate set yields a readable trajectory-based score. The frozen configuration uses seven candidate preferences
$\theta_k\in\{-3,-2,-1,0,1,2,3\}\times\pi/8$; the grid spans $[-3\pi/8,3\pi/8]$ inside the declared domain $[-\pi/2,\pi/2]$ of Eq.~\ref{eq:estimator}, and the admissible span quoted in Fig.~\ref{fig:monitor} is the candidate span $3\pi/4$ (Fig.~\ref{fig:monitor} caption). A rolling window of the ten most recent frames at
$10$\,Hz (one second), with at least four observed frames required for a reading; and a Gaussian
trajectory likelihood $\ell_{a,k}\propto\exp\{-\mathrm{MSE}_k/(2\sigma^2)\}$ with $\sigma=0.1$,
where $\mathrm{MSE}_k$ is the mean squared Euclidean distance (m$^2$) between observed and candidate
positions over the window, so $\sigma$ is in metres. The reported reading is not one of the seven
candidates: the likelihoods are normalised to weights $w_k=\ell_{a,k}/\sum_j\ell_{a,j}$ and the
reading is their weighted mean $\hat\theta_a=\sum_k w_k\theta_k$, which varies continuously over
$[-3\pi/8,3\pi/8]$. The candidate set is the support of the likelihood, not the range of the
estimator; the concentration of $w$ is what the readability gate below tests. A reading failing the near-uniform rule above carries no discriminative information between
the candidates, and the monitor abstains. The full frozen configuration is released with the code.

\subsection{Interaction opportunity, readability and episode summary}
An interaction-opportunity indicator $o_{ij}(t)$ marks when a behaviourally meaningful interaction is
active (a mapped conflict point reachable by both agents within a bounded horizon with closing
geometry), as opposed to mere co-presence. Candidate-weight concentration supplies the operational
candidate-discrimination statistic $u_a(t)$. The readability indicator is
\begin{equation}
  g_a(t) = \mathbf{1}\!\left[\,o_{ij}(t)=1,\ u_a(t)<u_0,\ q_a(t)>q_0,\ c_a(t)>c_0\,\right],
  \label{eq:readability}
\end{equation}
where $q_a(t)$ records trajectory/reference information and $c_a(t)$ counterpart-identity stability.
The thresholds are frozen on training/guard data and applied identically downstream; $g_a(t)=0$
states are labelled as no active interaction, insufficient candidate discrimination, unstable
identity or unreliable reference rather than scored.
The event-level summary of an interaction
is the mean current IPV over its valid interaction segment; alternative rules (all valid frames;
readability-weighted) differ by $0.26$ rad on average and flip the episode sign in $7$--$22\%$ of
episodes (Supplementary Fig.~\ref{fig:suppmeasurement}c), which is why the monitor never uses an episode summary
online and every reading is per-moment.

Three dyad-level analyses in Section 2.2 use the episode summary for both partners. The cohort is the
$34{,}850$ two-agent cases with a valid estimate for both partners. Complementarity treats the two
summaries as an exchangeable pair and compares their intraclass correlation with a null that re-pairs
partners across cases within $415$ scenario-matched strata (source subset, path-geometry category,
post-encroachment band and duration band); $34{,}757$ cases lie in strata that support re-pairing.
The estimate is computed in two independently built implementations and uncertainty is obtained by
scene-cluster bootstrap. Early legibility predicts, out of fold under grouped cross-validation, which
partner's summary will finish on the giving-way side from the first fraction of the interaction or
from early kinematics over the same window; discrimination is AUC with scene-cluster bootstrap
intervals. Turning yield is the paired difference between the turning and straight-driving partners,
reported by source and pooled with Benjamini--Hochberg correction across turn groups.

\subsection{Context-conditioned human reference range}
The reference range is the conditional distribution of the human current IPV given the current
observable situation $z_t$ (geometry, role, interaction progress, relative kinematics,
online-computable risk proxies, and support/readability state), estimated by conditional quantile regression (histogram-based gradient-boosted quantile models, one per quantile level, with frozen hyperparameters: learning rate $0.06$, $72$ boosting iterations, maximum $31$ leaf nodes, minimum $80$ samples per leaf, L2 regularisation $0.01$, $255$ bins, no early stopping)
and evaluated only on readable, active windows. The reference pool contains human--human vehicle
pairs only ($2{,}442{,}625$ anchor rows): any interaction in which either agent is an automated
data-collection vehicle is excluded from the reference, so machine behaviour cannot contaminate the
human range. The corpora are fleet-collected in specific cities in the United States and Singapore;
the reference is a sample of natural driving there, not of all human driving. We compare a global reference, an offline
oracle-risk ceiling, and the context-conditioned reference; the context-conditioned reference
supplies essentially all of the achievable sharpening over the global reference (width reductions of
21\%, 20\% and 8\% at the 80\%, 90\% and 95\% levels, at coverage within 0.6 percentage points of
nominal; the global 95\% range spans essentially the whole admissible scale). Two further inputs
are evaluated as ablations and add no measurable value: conditioning on the principal counterpart's
current IPV, and conditioning on the ego agent's own behavioural history (the self-anchor). Neither
is interpreted as a mechanism; the counterpart channel in particular is statistically indistinguishable
from removing the IPV input given the situation (paired $90\%$ interval-score difference $-0.0002$,
case-clustered $p=0.86$), which---set against the event-level partner
correlation---is the tension highlighted in the Results. Concretely, $z_t$ comprises four categorical descriptors (path-geometry category, path relation, turn-pair label and priority role) and $22$ numeric channels: elapsed interaction time and history length; ego and counterpart velocity components and headings at the anchor; relative position components, distance, velocity components, speed, closing rate and heading difference at the anchor; short-window means and a dispersion of relative distance, relative speed and closing rate; and two online risk proxies --- the closing-rate time-to-collision at the anchor and an anticipated post-encroachment proxy computed only from information available at time $t$.
Excluded from the reference are observed
post-encroachment time, realised passing order, post-hoc phase labels, the episode summary, the
machine/human identity, and target-proximal concurrent ego actions; quantiles are made non-crossing
by rearrangement. The reference range conditions on the current situation rather than on a specific
temporal pattern of the interaction; temporal-history features were examined and are not part of the
primary conditioning.

\subsection{Conformal calibration and abstention}
Vocabulary, used consistently: an anchor row is one agent-moment eligible for
estimation; a readable moment passes the discriminative-information rule; an accepted moment
additionally lies within human support (these receive verdicts); on the benchmark we write candidate
moment for an anchor row entering the two gates, and judgeable moment for one that passes both. A scene
is one contiguous recorded segment and the unit of the frozen split; an interaction case is one
interacting vehicle pair within a scene (Section 2.2); on the benchmark, a scenario is one of the staged
templates and a scenario run is one system's or one driver's traversal of it (Section 2.4).
The final interval is calibrated by split-conformal prediction~\cite{romano2019cqr,angelopoulos2023conformal} with nonconformity
$s_k=\max\{\hat q_l(x_k)-y_k,\ y_k-\hat q_u(x_k),\ 0\}$ and radius
$c_\alpha=s_{(\lceil(n+1)(1-\alpha)\rceil)}$ over accepted calibration moments.
Conformal calibration of trajectory atypicality has precedent in unconditional surveillance
settings~\cite{laxhammar2014online}; distribution-shift-aware conformal
variants~\cite{tibshirani2019conformal, gibbs2021adaptive, lei2018distribution} are compatible with
the same monitor, and we use the split-conformal form for auditability.
The estimand is pointwise marginal coverage over accepted moments: the conformal radius $c_\alpha$
is the $\lceil(n+1)(1-\alpha)\rceil$-th smallest nonconformity score over all accepted
calibration-fold moments ($n=455{,}723$), and coverage is evaluated over all accepted test-fold
moments ($461{,}937$). The fitted radii are near zero ($c_\alpha = 1.4\times10^{-3}$, $1.2\times10^{-6}$ and $0.0$\,rad at the $80\%$, $90\%$ and $95\%$ levels), so the conformal step finds essentially nothing to repair in the situational quantile model.
Splits are made by whole scenes under a split frozen before this study, so no
scene contributes to more than one fold; the readability gate is frozen on training/guard data and
applied identically to calibration and test. Moments within an interaction are dependent, so the
coverage statement is marginal over accepted moments --- not a per-interaction, sequential or
conditional guarantee --- and the achieved coverage is reported empirically (Fig.~\ref{fig:monitor}b). A marginal
statement leaves conditional coverage unconstrained in principle, so we also report it by source:
on the test fold the deployed reference contains $88.2\%$, $95.9\%$, $85.9\%$ and $88.2\%$ of
moments at the $90\%$ level within Waymo, nuPlan, Lyft and Argoverse-2 respectively, against $90.3\%$
pooled, so the pooled figure is not masking a source in which the range fails.
The frozen split assigns $974{,}984$ / $499{,}893$ / $481{,}088$ / $486{,}660$ human anchor rows to
the train / guard / calibration / test folds ($26{,}828$ interaction cases; whole-scene split frozen
in advance). The universes nest as follows: of the $38{,}228$ corpus interaction cases (Section 2.2),
$26{,}828$ enter the frozen anchor ledger, contributing $4{,}497{,}368$ anchor rows (the readability
accounting above); the human--human reference pool retains $2{,}442{,}625$ of those rows (Methods
4.3), which are exactly the train/guard/calibration/test folds
($974{,}984+499{,}893+481{,}088+486{,}660$). The support gate scores each moment by the mean Euclidean distance to its $k=25$
nearest training anchors after standardisation and one-hot encoding of the situation categories,
abstaining above the $95$th percentile of guard distances ($1.081$); categorical support
additionally requires at least $50$ training anchors and at least $10$ cases. On the test fold,
support abstention is $5.08\%$ ($24{,}723$ of $486{,}660$) with zero categorical failures. On the
natural corpus the mechanism-level readability accounting over $4{,}497{,}368$ anchor rows is:
$3{,}202{,}646$ readable, $1{,}275{,}480$ near-uniform weights, $17{,}416$ exact ties and $1{,}826$
solver failures. The $1{,}826$ solver failures belong to this natural-corpus accounting; on the benchmark, no
candidate moment failed the solver or pipeline (Section 2.4). Conditional coverage across state groups and trajectory-wise simultaneous coverage
are not established here and remain stated boundaries of the construction; leave-one-source-out
coverage quantifies the source-shift boundary: at the $90\%$ level, refitting with one source held out
covers the held-out source at
$0.743$ (Waymo, $143{,}380$/$193{,}096$), $0.990$ (nuPlan, $149{,}068$/$150{,}587$), $0.750$ (Lyft,
$68{,}270$/$91{,}069$) and $0.900$ (Argoverse-2, $11{,}315$/$12{,}576$, with $44.3\%$ held-out
abstention---the support gate intercepting unsupported situations); transfer across sources is
therefore a bounded property of source heterogeneity, and coverage and abstention must be read
together. We report not only accepted-window coverage but also the active-interaction rate,
readability rates, and the total abstention rate with its reason distribution, so that filtering is
never hidden. The monitor returns \textsc{Abstain}---
not \textsc{Within-Norm}---with one of the reason codes \textsc{No-Active-Interaction},
\textsc{Ego-Unreadable}, \textsc{Counterpart-Unreadable} (the counterpart's identity is not stable
enough to anchor a reading, $c_a$ of Eq.~\ref{eq:readability}), \textsc{Map-or-Role-Uncertain},
\textsc{Low-Human-Support} or \textsc{Out-of-Distribution}; reason codes are assigned by the first failing gate in the evaluation order of Algorithm 2, so exactly one code is returned; otherwise it emits \textsc{Within-Norm},
\textsc{Competitive-Deviation} or \textsc{Over-Yielding}. In the main text,
\textsc{Competitive-Deviation} is written `outside the range on the assertive side' and
\textsc{Over-Yielding} `outside on the accommodating side'; both sides are tested at every
judgeable moment. A sequential warning layer (persistence over consecutive verdicts) belongs to the
deployment interface: no result in this paper uses it, every reported verdict and flag rate is
per-moment, and its operating characteristics (episode-level false-alarm rate, detection delay,
warning duration) are not evaluated in this paper.

\subsection{Real-vehicle consequence protocol}
We use a matched-scenario real-vehicle benchmark at the scenario level (no run-level, seed-level or
causal claims), separating an outcome universe from a replay/IPV universe with frozen exclusions.
The counterpart vehicles are driven by traffic-simulation software (TESS~NG) and respond to what the
ego vehicle does; the same counterpart control is used in the human reference arm below, so the two
arms face a counterpart that reacts to them on the same terms. The staging is mixed-reality: the ego
is a real vehicle driven at a test site, and the simulated counterparts reach an automated system
through its perception interface and a human driver through a head-mounted display. No
confirmatory analysis proceeded until an outcome-blind crosswalk had reconciled canonical scenario
identifiers, structural case identities and official outcome labels. Realised-consequence events are
defined by a frozen automatic extractor. Verdicts are issued per moment; outcome windows begin at
the flagged moment, so the trajectory samples that produce a verdict (the one-second window ending
at that moment) precede every outcome sample. The ego-side outcome is the minimum time-to-collision
to the counterpart over the fixed three-second post-verdict window ($30$ frame intervals at
$10$\,Hz, inclusive of both endpoints; the run ends inside that window for $4.8\%$ of assertive,
$6.2\%$ of accommodating and $5.0\%$ of within-range moments, and the window is then truncated at
the last available frame); time-to-collision is computed per frame as distance over closing rate,
and frames in which the pair is not closing do not enter the minimum; the counterpart-side outcomes
are read over the same fixed three-second window from the same moment, from the other vehicle's own
logged control record.
Comparisons are made within the frozen situation cells, inference resamples scenario runs (the unit
of clustering throughout), and a placebo test that reassigns whole exposure trajectories across
scenario runs distinguishes the real flag timing from permuted timings (200 whole-trajectory reassignment draws; the comparison statistic is the absolute case-clustered $t$; empirical $p=0.0199$).
Across all six level--window combinations of the $\log(1+\mathrm{TTC})$ magnitude battery the
coefficients agree in sign; the fixed three-second window is the pre-specified primary and is the one
displayed, and its case-clustered intervals exclude zero at all three levels, while the alternative
and much shorter window---the anchor's own prediction-target horizon, a median $0.60$\,s---crosses
zero at the $90\%$ level ($[-2.6100, +0.1372]$). A case-level label permutation
on the same battery does not reach significance ($p=0.1493$); the exposure placebo above is the
test specific to flag timing.
The battery is reported for both sides of the range. Its endpoints, windows, thresholds, clustering
unit and resampling settings were fixed for the assertive-side comparison and were not altered
afterwards; the accommodating-side values were obtained by running that unchanged battery on the
accommodating group, so no endpoint, threshold or window was chosen in the light of
accommodating-side data. The counterpart braking endpoint was specified for both sides from the
outset. On the accommodating side ($817$ ego moments and $719$ counterpart anchors, against the same
within-range groups) the counterpart's speed reduction is $0.79\times$ the within-range median
($[0.28,1.53]$) and its speed variation $1.08\times$ ($[0.78,1.53]$); the share of counterpart
records braking below $-3$\,m\,s$^{-2}$ differs by $+0.07$ percentage points ($[-1.96,+2.24]$); and
the ego's own margin median is $0.99\times$ ($[0.83,1.22]$) with its upper quartile at $1.03\times$
($[0.67,1.38]$). Every interval admits no difference.
These are non-detections at the achieved precision---which on the ratio scale is slightly coarser
than on the assertive side, despite the larger group---and not demonstrations of equivalence.
The two sides are additionally compared with each other, rather than each against the within-range
group in turn: on every resampled set of scenario runs both ratios are recomputed and their
difference taken, so the two sides move together and the comparison keeps the correlation that
separate intervals discard. At the upper quartile the sides differ ($-0.43$ $[-0.71,-0.06]$,
$p=0.014$); at the median they do not ($-0.09$ $[-0.29,+0.12]$, $p=0.368$). That ordering is
unchanged on the shorter window and under the retention convention above---the upper-quartile
difference excludes zero in all four window-by-convention combinations and the median difference in
none---so the side-specific statement is made at the quartile that supports it and withheld at the
one that does not.
Estimates are weighted by moment rather than by scenario run, because the monitor issues a verdict
at each moment and the per-moment population is the one its operation defines; resampling scenario
runs sets the intervals. Weighted equally per run, the assertive-side counterpart braking contrast
does not reach significance ($p=0.4704$); the speed-reduction and speed-variation contrasts carry
the counterpart-side result. Ratios of medians are reported alongside differences because a ratio
becomes unstable where the within-range median is near zero, as it is in one of the two sites.
Of the $14{,}099$ judgeable
moments at the $90\%$ level, $519$ lie below the range, $869$ above and $12{,}711$ inside.
The criterion is the situation-conditioned range, not the sign of the reading: at the $90\%$ level the
lower edge is negative in essentially every situation (median $-1.03$\,rad), $40.0\%$ of within-range
moments also carry a negative reading, and of all negative-reading moments only $9.25\%$ are flagged.
The challenge fielded $20$ systems; one is set aside in full because its replay records are not
clean, leaving $19$ systems across the $15$ scenarios, or $285$ system--scenario cells, of which
$267$ yield the trajectory anchors the estimator needs. The $18$ that do not were all attempted and
officially scored ($15$ scored completions and $3$ collision failures); they are lost in our own
replay processing --- $11$ to observed reference trajectories that collapse to fewer than two
distinct points, $3$ to no timing-valid anchor, $3$ to no eligible counterpart and $1$ to a
counterpart without timing support --- and $13$ of the $18$ fall in a single system. No scenario is
dropped: each retains between $16$ and $19$ systems.
Consequence analyses follow the same run accounting as the rest of the benchmark. Of the $267$
scenario runs, $231$ contain at least one judgeable moment and can therefore enter a consequence
analysis at all. Reported per run rather than per moment, $157$ of those $231$ ($68.0\%$) contain at
least one flagged moment on either side of the range and $120$ ($51.9\%$) at least one on the
assertive side. The per-moment rate is the operating point throughout; the per-run figure is given
here because the two have different denominators and are easily confused. Of those, $27$ are runs in which the counterpart was identified from a different
naming convention in the run logs; they are held out of the pooled counterpart-response analyses and
reported separately, leaving $204$, of which $175$ have the recorded counterpart logs the outcome
battery requires. The split is one of record-keeping rather than of apparatus: the counterpart is
the same kind of simulator-driven vehicle throughout. The runs dropped at that last step do differ
from the retained ones in how often they contain a flagged moment: $7$ of the $29$ dropped ($24.1\%$)
against $132$ of the $175$ retained ($75.4\%$). The requirement is a logging one, and a run without
recorded counterpart logs also tends to be one with fewer judgeable moments, so the two rates are
not independent of each other; we report them rather than argue the point. The restriction applies
only to the counterpart-side panels: the ego-side panels use all $227$ runs with an evaluated ego
window and are not filtered on counterpart logging at all. That the counterparts respond is also visible in the replay itself: in the four
scenarios whose counterpart identity recurs across systems, the same counterpart's trajectory
differs between systems by $4.35$--$9.03$\,m, and in none of the four is it identical. The ego-margin panels use the moments with
a defined post-verdict margin ($499$ assertive, $817$ accommodating, $12{,}344$ within-range). The
margin is undefined where the pair is not closing, and the three groups do not lose equally: $3.9\%$
of assertive moments ($20$ of $519$), $6.0\%$ of accommodating moments ($52$ of $869$) and $2.9\%$
of within-range moments ($367$ of $12{,}711$). The accommodating group loses most, which follows
from what defines it---yielding is what stops a pair closing---so the exclusion removes
preferentially the moments in which that group is most itself, and it does so in the direction that
draws the accommodating group toward the within-range one. We therefore report the battery again
with those moments retained and ranked at the safe end rather than dropped, which is what never
approaching a collision means; the quartile statistics are invariant to the value used, verified
against two sentinels a thousandfold apart. Retaining them leaves the assertive
upper-quartile contrast in place ($0.626$ of the within-range quartile, $[0.410, 0.829]$, against
$0.601$ $[0.436, 0.831]$ when they are dropped) and moves the accommodating side further away from
the within-range group in the direction of wider margins ($1.029$ to $1.126$), which is what
restoring the most strongly yielding moments should do. The exclusion was therefore pulling the
accommodating side toward the within-range group, not the assertive side. Under these definitions the harm
relationship is bounded rather than confirmed: deviation does not robustly or specifically predict
realised harm, and the dominant limitation is upstream, in the adequacy of interaction-failure
segmentation itself. The fraction of moments whose reading carries discriminative information
differs between the naturalistic corpora and the benchmark ($70.3\%$ vs $55.3\%$); we report this as
an observed difference without attributing a cause, and the readability gate treats both populations
identically. The binding constraint on benchmark coverage is situational rather than volumetric: in the 36 always-silent runs the readability pass rate ($53.78\%$) is close to the benchmark average ($55.30\%$) while the human-support pass rate is $0.07\%$ (against $32.32\%$ overall), and support does not follow sample volume --- one high-volume situation cell with $1{,}148{,}133$ human anchor rows passes support on $14.58\%$ of its benchmark moments while a $45{,}283$-row cell passes on $47.03\%$. The support gate compares kinematic neighbourhoods, not IPV values.

\subsection{Human reference arm (matched scenarios)}
Licensed drivers drove the same scenario set in a real vehicle against the same simulator-driven
counterparts under the counterpart control described above,
giving a drivers-by-scenarios matrix (20 drivers $\times$
15 scenarios) symmetric to the automated-systems benchmark. The unchanged estimator,
abstention gates and frozen reference range were applied to the human runs. Endpoints were fixed
before analysis: (i) the flagged fraction of judgeable moments against the nominal level, with the
natural-driving outside rate ($9.72\%$ at the $90\%$ level) as the apparatus-validity anchor; and
(ii) the consequence battery of Fig.~\ref{fig:consequence} recomputed within the human arm under
identical frozen definitions, with no new outcome variables and no threshold scans. Official
competition scores, harm labels and preference ratings were excluded as endpoints in the analysis
plan, which was specified and frozen before any outcome was examined; those labels are publicly
released by the challenge, which is operated by the authors' research group, and no author entered
the challenge as a competitor. Driver identities were pseudonymised before analysis.

The human-reference study was approved by the Science and Technology Ethics Committee of Tongji
University (approval tjdxsr2025011). All participants gave written informed consent before taking
part and were insured for the duration of testing. Runs took place at a dedicated test site with a
safety supervisor on board who could intervene or terminate a run at any time. Participants held a
valid driving licence, were compensated at CNY 150 per hour, and were instructed to drive naturally;
they were not told that their runs would serve as a human reference against automated systems. Each
driver completed the scenario set as a single fixed sequence, identical across drivers, and no trial
was missing.

\subsection{Subjective evaluation of selected interaction segments}
\label{sec:subjective_methods}
\noindent\textbf{Cohort and stimuli.} The author-confirmed final rating dataset contains 40
participants organised into 20 pairs and 40 sessions. Of 600 formal trial attempts, 596 were completed
and four were aborted. Each completed trial contributes one ego-position and one counterpart-position
rating record, yielding 1{,}192 records; repeated records from the same participant are not treated
as independent participants. The 90 segments span 15 scenarios: 24 assertive-side segments,
36 accommodating-side segments and 30 within-range comparison segments, contributing 158, 239 and
199 completed trials, respectively. The supplied source labels identify 45 automated and 45 human
segments. These labels are used for sensitivity summaries, not a source-recognition or vehicle-ranking
claim. The subjective protocol is distinct from the 20-driver matched driving arm.

The supplied segment categories were retained. Package-internal checks associate A with segments
containing below-range moments, C with segments containing above-range moments, and W with
within-range and abstention moments. A segment category does not assign the same verdict to every
frame; abstention is never recoded as a within-range verdict. Stimulus-selection rules and the
version relationship between these supplied labels and the manuscript's frozen monitor require
completion in the study documentation below.

\noindent\textbf{Rating items.} Analyses are separated by the supplied ego and counterpart position
labels. Perceived atypicality is Q1 reverse-coded as $8-\mathrm{Q1}$; Q4 is comfort, with higher
values indicating greater comfort; Q5 is the supplied reported interaction-cost item, with higher
values indicating greater cost. These items use a 1--7 scale. Q5 and Q5b are structurally unavailable
for the ego position: the 596 blank values for each item there are not zero ratings. The four
aborted trials contribute no completed questionnaires and are not imputed. Supplementary
Note 2 records the remaining exploratory items and the scope of their interpretation.

\noindent\textbf{Statistical analysis.} This is an exploratory secondary analysis of the final
rating cohort. Within each participant and position, ratings were averaged within category; category
differences were then computed within participant and averaged with equal participant weights.
Pointwise $95\%$ percentile intervals use 10{,}000 bootstrap resamples of the 20 participant pairs,
keeping the two members together. Two-sided one-sample $t$-tests use the 20 pair-averaged differences
($19$ degrees of freedom). Holm correction covers all 30 available item-by-position-by-category
contrasts, including contrasts not highlighted in the main text. Bootstrap intervals are pointwise,
not multiplicity-adjusted simultaneous intervals, and condition on the selected stimulus set.
The extraction used Python 3.9.6, NumPy 2.0.2, pandas 2.3.3 and SciPy 1.13.1, with seed 20260925.

A separate scenario-composition sensitivity analysis first averages ratings within segment, then
segments within scenario and category, and finally contrasts the categories across shared scenarios
with equal scenario weights. Assertive-minus-within, accommodating-minus-within and
assertive-minus-accommodating comparisons contain eight, nine and 12 shared scenarios, respectively.
Its intervals resample scenarios and condition on the observed raters; they are not a joint
participant-and-stimulus uncertainty estimate. The participant-pair and scenario analyses address
different sources of variation and are not independent replications.

\planned{Complete the subjective study's recruitment and eligibility, apparatus and presentation
instructions, allocation and order rules, compensation, ethics approval and consent scope, full
questionnaire wording and anchors (especially Q5), stimulus-selection rules, and the version link
between supplied labels and the frozen monitor. These records were absent from the extraction
package; do not substitute the matched human-driving protocol.}

% ---------------------------------------------------------------------
\begin{algorithm}[t]
\caption{Offline calibration of the context-conditioned reference range}
\label{alg:calib}
\begin{algorithmic}[1]
\REQUIRE existing IPV time series $\mathcal{D}$; target coverage $1-\alpha$
\ENSURE quantile model $\hat{Q}$, conformal radius $c_\alpha$, frozen readability gate
\STATE split whole scenes into train / guard / calibration / test
\STATE freeze the interaction-opportunity rule and the readability gate ($u_0,q_0,c_0$) on train/guard
\FOR{each interaction in train}
    \STATE keep active, readable windows; build the current situation $z_t$; target $y\leftarrow$ human current IPV
\ENDFOR
\STATE fit non-crossing situational quantiles; fit counterpart-IPV and self-history variants as ablations only
\STATE apply the identical frozen gate to calibration; score every accepted calibration window
\STATE $c_\alpha\leftarrow s_{(\lceil(n+1)(1-\alpha)\rceil)}$ over accepted calibration windows; freeze all parameters
\RETURN $\hat{Q},\,c_\alpha,\,\text{gate}$ \COMMENT{evaluate once on test}
\end{algorithmic}
\end{algorithm}

% ---------------------------------------------------------------------
\begin{algorithm}[t]
\caption{Online social-atypicality monitoring step at time $t$}
\label{alg:online}
\begin{algorithmic}[1]
\REQUIRE window $W_t$; map-lane reference; calibrated $\hat{Q},c_\alpha$; frozen gate; source health
\ENSURE reference range and verdict, or an abstention with reason code
\STATE identify the principal counterpart; evaluate active interaction $o_{ij}(t)$
\IF{$o_{ij}(t)=0$} \RETURN \textsc{Abstain}(\textsc{No-Active-Interaction}) \ENDIF
\STATE estimate $\hat\theta_i(t)$ and reliability $u_i(t)$ \COMMENT{Eq.~\eqref{eq:estimator}}
\IF{$g_i(t)=0$} \RETURN \textsc{Abstain}(\textsc{Ego-Unreadable} if $u_i$ or $q_i$ failed; \textsc{Counterpart-Unreadable} if $c_i$ failed) \ENDIF
\STATE build the current situation $z_t$; run map/role, human-support and distribution gates
\IF{any gate fails} \RETURN \textsc{Abstain}(reason code) \ENDIF
\STATE $[\hat q_l,\hat q_u]\leftarrow \hat{Q}(z_t)$ widened by $c_\alpha$
\STATE verdict $\leftarrow$ \textsc{Within-Norm} / \textsc{Competitive-Deviation} / \textsc{Over-Yielding} from $\hat\theta_i(t)$ vs $[\hat q_l,\hat q_u]$
\STATE optionally map the verdict to a planner-facing signal through a persistence rule (deployment interface; not evaluated in this paper)
\RETURN (range, verdict) \COMMENT{plus signal only where the optional persistence rule is instantiated}
\end{algorithmic}
\end{algorithm}

\noindent\textbf{Planner interface.} The verdicts map to planner actions---audit,
soft cost or warning, a fallback candidate under sustained competitive deviation at high physical
risk, a hesitation/efficiency warning, and --- on abstention --- a degraded mode in which the social verdict is simply unavailable. This mapping is an interface specification, not an external-validity result; the social layer is advisory and subordinate to collision-safety enforcement, no social verdict or abstention relaxes or delays any safety constraint, and a planner fallback is a control action never labelled \textsc{Within-Norm}.

\section*{Data Availability}
InterHub-derived interaction cases and source datasets (Waymo, nuPlan, Lyft, Argoverse-2) are governed
by their respective licences; processed derivatives, including the online IPV time series, will be made
available upon publication. The matched-scenario real-vehicle benchmark is operated by the authors'
research group at Tongji University, and its official outcome labels are publicly released; the
derived moment-level verdict series for the automated-systems arm will be made available upon
publication. Records from the human reference arm are not released: participants consented to take
part in the study, and that consent does not extend to sharing or publishing their driving records,
so the human arm is reported only in aggregate.

\planned{Specify the separate subjective study's data-access and sharing terms from its own
consent records, including what may be released as aggregate Source Data and whether any
participant-level derivatives may be shared. The driving-arm restriction above does not establish
the subjective study's permissions.}

\section*{Code Availability}
Implementation, the online likelihood-based inverse-planning IPV estimator, the readability gate and
abstention logic, and figure-generation code will be deposited in Zenodo under a persistent
identifier upon publication, and are available to referees on request during peer review.


% ======================================================================
% Supplementary Information
% ======================================================================
\clearpage
\section*{Supplementary Information}
\setcounter{figure}{0}
\renewcommand{\thefigure}{S\arabic{figure}}
\setcounter{table}{0}
\renewcommand{\thetable}{S\arabic{table}}

\subsection*{Supplementary Note 1. Published validation and present implementation of IPV estimation}
The IPV estimator used in this study was introduced in Zhao et al.~\cite{zhao2024measuring}. In
controlled VGIM interactions, agents were assigned interaction-preference patterns and the
trajectory-based estimate tracked those imposed patterns as the interaction unfolded; incorporating
the identified IPV also improved interactive-motion prediction. The present study adopts that
published estimator without refitting it and changes the downstream question from parameter
identification to membership in a situation-conditioned human reference range.

The online implementation evaluates seven candidate preferences
$\{-3,-2,-1,0,1,2,3\}\times\pi/8$ over a one-second causal window at $10$\,Hz. Each candidate
generates a trajectory under the frozen VGIM utility, the observed-to-candidate trajectory error is
converted to a normalised likelihood, and Eq.~\ref{eq:estimator} returns their continuous weighted
mean. A static route-conditioned map reference keeps the calculation causal and available at decision
time. The operational readability gate asks whether the current candidate set is sufficiently
discriminated while the interaction, trajectory reference and counterpart identity are valid.

The present implementation was also checked against interaction-specific controls
(Supplementary Fig.~\ref{fig:suppmeasurement}). Candidate discrimination sharpens in the actual dyad
relative to time-misaligned, partner-swapped and no-opportunity controls. At the same time, readable
frame-level values can continue to move as negotiation develops, and episode-level summaries depend
on how valid frames are aggregated. These diagnostics support the per-moment design used throughout
the paper: the monitor retains the frame-level reading, applies an explicit readability boundary and
never substitutes a single agent-level score for the online trajectory.

\begin{table}[!htb]
\centering
\small
\caption{\textbf{Role of IPV estimation in the published T-ITS study and in the present monitor.}}
\label{tab:suppipv}
\begin{tabular}{p{0.20\linewidth}p{0.35\linewidth}p{0.35\linewidth}}
\toprule
Aspect & Published T-ITS study~\cite{zhao2024measuring} & Present study \\
\midrule
Scientific role & Introduces trajectory-based IPV identification in the VGIM framework & Uses the published estimator as a fixed input to a human-reference monitor \\
Interaction setting & Controlled pairwise VGIM interactions & Naturalistic and matched-scenario pairwise interactions \\
Output & Time-varying IPV identified from observed interaction trajectories & Frame-level IPV reading compared with the human range for the current situation \\
Primary question & Can imposed interaction-preference patterns be identified and improve interaction prediction? & Does the current behaviour fall within the context-conditioned human reference range? \\
Validation emphasis & Model-consistent identification and interactive-motion prediction & Context conditioning, calibration, abstention and behavioural auditing \\
\bottomrule
\end{tabular}
\end{table}

\begin{figure}[!htb]\centering
\includegraphics[width=\linewidth]{figures/fig1_measurable.pdf}
\caption{\textbf{Interaction-specific diagnostics and summary sensitivity of the present online implementation.}
\textbf{a} Change in candidate-weight concentration inside real interaction windows relative to three
controls. The interaction-specific sharpening disappears under time misalignment ($+0.006$), is
weaker under partner swapping ($-0.043$), and reverses where no interaction opportunity exists
($+0.027$). Points are means with $95\%$ confidence intervals over case clusters; the $n$ column
counts case clusters.
\textbf{b} Mean movement of the frame-level reading between consecutive frames. Readable values remain
dynamic as an interaction unfolds, so the monitor keeps them time-indexed.
\textbf{c} Episode-level summaries of the same interactions under three aggregation rules; bars span
the interquartile range and ticks mark medians. The rules differ by $0.26$\,rad on average and change
the sign in $7$--$22\%$ of episodes, motivating the per-moment verdict used in the main analysis.
Source data are provided as a Source Data file.}
\label{fig:suppmeasurement}
\end{figure}

\subsection*{Supplementary Note 2. Subjective endpoints and exploratory comparisons}
\label{supp:subjective}
The main subjective results focus on perceived atypicality, comfort and reported counterpart
interaction cost. All 30 item-by-position-by-category comparisons were retained in the extraction,
including the supplied Q2 yielding-tendency item, Q3 predictability item and counterpart-only Q5b
naturalness item. Their full wording and response anchors must accompany the final questionnaire.
No item is substituted for a direct preference, acceptability or appropriateness question.

The ego-position assertive-minus-within atypicality difference is $+0.465$
$[+0.170,+0.759]$ points, with Holm-adjusted $p=0.0704$. It is therefore not described as a
multiplicity-adjusted positive finding despite its pointwise interval excluding zero. The
counterpart-position assertive-minus-accommodating atypicality comparison has adjusted $p=1$;
this non-detection does not establish equal atypicality. Q2 does not clearly separate assertive
from within-range segments, and the Q5b comparisons do not establish equivalence of naturalness.
Continuous dial responses are not used as evidence of instantaneous discomfort or warning latency:
their construct, anchors, standardisation and onset definitions require additional documentation.

Table~\ref{tab:subjective_core} gives the eight counterpart-position contrasts highlighted in the
Results. Participant-pair intervals condition on the selected stimuli. The common-scenario estimate
uses equal weighting of segments within scenario and then of shared scenarios; it is a composition
sensitivity summary with a different estimand, not a second experiment.

\begin{table}[!htb]
\centering\small
\caption{\textbf{Principal counterpart-position subjective contrasts.}
Differences are on 1--7 scales; positive atypicality and cost differences indicate higher ratings,
while negative comfort differences indicate lower comfort. A, assertive-side; C, accommodating-side;
W, within-range comparison. All rows use 40 participants in 20 pairs. Intervals are pointwise
$95\%$ participant-pair bootstrap intervals. Adjusted $p$ values use the full 30-comparison Holm
family. Common-scenario means use eight A--W, nine C--W and 12 A--C scenarios.}
\label{tab:subjective_core}
\setlength{\tabcolsep}{4pt}
\begin{tabular}{llrrr}
\toprule
Item & Contrast & Difference [$95\%$ CI] & Holm $p$ & Scenario mean \\
\midrule
Atypicality & A--W & $+0.639\ [ +0.318,+0.992]$ & $0.0245$ & $+0.850$ \\
Atypicality & C--W & $+0.682\ [ +0.548,+0.817]$ & $1.99\!\times\!10^{-7}$ & $+0.949$ \\
Comfort & A--W & $-1.091\ [-1.226,-0.961]$ & $6.11\!\times\!10^{-11}$ & $-1.120$ \\
Comfort & A--C & $-0.870\ [-0.974,-0.766]$ & $5.18\!\times\!10^{-11}$ & $-0.851$ \\
Comfort & C--W & $-0.221\ [-0.351,-0.100]$ & $0.0342$ & $-0.245$ \\
Reported cost & A--W & $+2.070\ [ +1.885,+2.263]$ & $3.95\!\times\!10^{-13}$ & $+2.005$ \\
Reported cost & A--C & $+1.163\ [ +1.019,+1.300]$ & $6.93\!\times\!10^{-11}$ & $+1.191$ \\
Reported cost & C--W & $+0.908\ [ +0.771,+1.051]$ & $3.49\!\times\!10^{-9}$ & $+0.915$ \\
\bottomrule
\end{tabular}
\end{table}

% ======================================================================
% Extended Data figures
% ======================================================================
\clearpage
\section*{Extended Data}
\suppressfloats[t]
\setcounter{figure}{0}
\renewcommand{\thefigure}{ED\arabic{figure}}

\begin{figure}[!htb]\centering
\includegraphics[width=\linewidth]{figures/figE2_reference_band.pdf}
\caption{\textbf{The human reference range rises, falls and changes width
with the situation.} A single interaction segment showing the frame-level reading against the nested
$80\%$, $90\%$ and $95\%$ situation-conditioned human reference ranges. Grey intervals in the top strip
mark frames with no reading; those frames remain gaps rather than being plotted as a value. Readable
frames are shown as inside the $90\%$ range, below it (more assertive, atypical) or above it
(the accommodating side, atypical). The two outlined frames carry the same reading, $-0.393$\,rad: it lies inside the range
where that range spans $2.19$\,rad and outside it where the range spans $1.49$\,rad. The pair isolates
the role of conditioning, in that the number being judged is identical while the reference against
which it is judged is not. The display is descriptive. Source data are provided as a Source Data file.}
\label{fig:edband}
\end{figure}

\begin{figure}[!htb]\centering
\includegraphics[width=\linewidth]{figures/figE1_case_example.pdf}
\caption{\textbf{One illustrative run containing flagged assertive moments.}
\textbf{a} Bird's-eye trajectories over the displayed window; red points mark the monitored vehicle's
positions during the selected flagged stretch.
\textbf{b} Speed timelines over a $2.0$\,s pre-flag baseline and the flagged stretch. Lines are centred
$21$-frame medians of speed computed from the logged velocity components. The counterpart median moves
from $4.31$ to $1.40$\,m\,s$^{-1}$; the automated vehicle's median moves from $3.33$ to
$3.32$\,m\,s$^{-1}$.
\textbf{c} The same stretch relative to the $90\%$ situation-conditioned human reference range. Two
traces are drawn from the same data: the faint line and its markers are the per-frame readings, which
are what the verdicts are computed on, and the bold line is those readings under a centred
$2.1$-second median---the filter the speed panel above also uses---which is what a display would
show. Filtering is applied to the shown value only and never to a verdict; the bold line is masked to
readable frames so it does not bridge an abstention. Red downward triangles are per-frame readings
below the lower edge of the range; grey ticks mark no reading, and available rows are connected only
for display continuity.
This run was selected from the $120$ runs containing at least one flagged assertive moment, and is one
of only two meeting display criteria fixed in advance (at least five flagged frames in one cluster,
successive flags separated by at most one second; counterpart speed decrease of at least $20\%$;
automated-vehicle speed change within $10\%$).
Of the two runs meeting all criteria, the displayed one ranks first under the selection score fixed
together with them --- the counterpart's speed-drop fraction minus the automated vehicle's absolute
speed-change fraction ($0.675$ vs $0.552$).
It is one illustrative run: it is not a typical value, not an overall distribution, and not evidence
for ranking or comparison between automated systems. The display is descriptive and assigns no
causation. Source data are provided as a Source Data file.}
\label{fig:edcase}
\end{figure}

\clearpage
\bibliographystyle{unsrt}
\bibliography{bibliography/biblio}

\section*{Acknowledgements}
This work was supported by the National Natural Science Foundation of China (Key Program, grant
52232015). We thank the drivers in the human-reference study and the participants in the subjective evaluation.

\section*{Author Contributions}
X.Z. conceived the study, developed the methodology and the software, and wrote the original draft
and the revisions. J.S. acquired funding and revised the manuscript. X.J. contributed validation and
data curation. Y.L. contributed methodology, data curation and validation. G.M. contributed to the
conception of the study and revised the manuscript.

\section*{Competing Interests}
The matched-scenario real-vehicle benchmark analysed in this work is organised by the authors'
research group at Tongji University. No author entered the challenge as a competitor, the official
outcome labels examined here are publicly released, and the funding for this study is independent of,
and does not overlap with, the funding of the challenge series. The authors declare no other
competing interests.

\end{document}
